\section{Approach} \label{sec:solution}
OWL-SDA~(OWL Synthetic Data AI-Agent) is an open-source\footnote{\url{https://github.com/milieuinfo/owl-sda}} Java toolkit for generating synthetic RDF data from ontology context and user constraints\footnote{\url{https://purl.org/owl-sda}}. Early design prototypes of the toolkit showed two recurring problems: direct Turtle editing by agents produced syntax errors and duplicate triples, while delayed validation caused an accumulation of violations. OWL-SDA addresses both by moving edits into an in-memory triple store and by exposing SHACL validation as a continuous feedback loop, which helps mitigate the incompleteness and inaccuracy risks also identified by Hao~et~al.~\cite{hao2024synthetic} in the creation of synthetic data.

The final design follows role-based task decomposition~\cite{NEURIPS2023_a3621ee9} and hierarchical supervision patterns~\cite{banihashemi2018hierarchical,long2026ai}. Workers handle narrow shape-level tasks that are sorted based on their complexity and links to other shapes, the supervisor coordinates scope and consistency (e.g.,~URI consistency or scenario), and the reviewer evaluates the final graph from a read-only perspective. This separation matters because SHACL-valid output can still contain invented vocabulary, inconsistent naming or implausible combinations across shapes.

\begin{figure}
    \centering
    \includegraphics[width=0.95\linewidth]{images/architecture.pdf}
    \caption{Architecture of the OWL-SDA pipeline}
    \label{fig:architecture}
\end{figure}

Figure~\ref{fig:architecture} summarises the workflow. During preparation, OWL-SDA loads the ontology, downloads imported vocabularies and derives two SHACL views: one used by workers for strict conformance checks and one used later by the reviewer for broader assessment. The supervisor then assigns shape-level tasks to workers, who add and remove triples in a shared concurrent store rather than editing Turtle text directly. The store performs duplicate checks and supports immediate syntax validation, allowing workers to repair local violations before moving on.
In a typical run, the supervisor reads the ontology, SHACL shapes and user context, workers generate triples for a small set of target shapes, each worker validates and repairs its own output, the supervisor validates the combined graph and dispatches targeted fixes for cross-shape inconsistencies, and finally a reviewer inspects the serialised output for broader plausibility issues. This keeps workers focused on local conformance while reserving global quality control for the more capable roles. OWL-SDA currently supports GitHub Copilot, Ollama and OpenAI-compatible endpoints, with provider selection configurable globally or per role. In practice, we used smaller, faster models for workers and stronger models for supervision and review, which follows cost-quality trade-offs also observed in multi-agent systems~\cite{chen2023less}.

During its development, we tweaked the workflow to introduce additional guardrails to catch common errors. For example, workers would sometimes create instance data that was using a different URI schema than the ontology, which was not caught by SHACL validation. Additional checks were added as system prompts to the supervisor and reviewer to help guide workers towards producing more consistent output.